% Standalone resolution manuscript generated from solution.md.
% Compile with XeLaTeX; author and review metadata are maintained in Markdown.
\documentclass[10pt,a4paper]{article}
\usepackage[margin=24mm,headheight=15pt]{geometry}
\usepackage{fontspec}
\setmainfont{DejaVuSerif.ttf}[BoldFont=DejaVuSerif-Bold.ttf,ItalicFont=DejaVuSerif-Italic.ttf,BoldItalicFont=DejaVuSerif-BoldItalic.ttf]
\setsansfont{DejaVuSans.ttf}[BoldFont=DejaVuSans-Bold.ttf,ItalicFont=DejaVuSans-Oblique.ttf]
\setmonofont{DejaVuSansMono.ttf}
\usepackage{amsmath,amssymb,unicode-math}
\setmathfont{latinmodern-math.otf}
\usepackage{xcolor,microtype,parskip,needspace,fancyhdr}
\usepackage{bookmark}
\usepackage{xurl}
\hypersetup{colorlinks=true,urlcolor=blue,linkcolor=blue,pdftitle={IE-26:
Sharp stability bounds for perturbed Fourier
interpolation},pdfauthor={George Stepaniants},pdflang={en-GB}}
\urlstyle{same}
\setlength{\emergencystretch}{3em}
\providecommand{\tightlist}{\setlength{\itemsep}{0pt}\setlength{\parskip}{0pt}}
\providecommand{\pandocbounded}[1]{#1}
\setcounter{secnumdepth}{0}
\allowdisplaybreaks[1]
\widowpenalty=10000
\clubpenalty=10000
\pagestyle{fancy}
\fancyhf{}
\fancyhead[L]{\small\sffamily AFFIRMATIVE RESOLUTION / IE-26}
\fancyhead[R]{\small\sffamily 12 September 2026 (UTC)}
\fancyfoot[L]{\parbox[b]{0.9\textwidth}{\fontsize{8}{10}\selectfont Substantial
AI assistance; independent automated review is documented separately.}}
\fancyfoot[R]{\footnotesize\thepage}
\begin{document}
{\raggedright\Large\bfseries IE-26: Sharp stability bounds for perturbed
Fourier interpolation\par}
\vspace{0.4em}
{\normalsize\bfseries George Stepaniants\par}
{\small Department of Computing and Mathematical Sciences, California
Institute of Technology, Pasadena, California, USA\par}
\vspace{0.6em}
This manuscript proves the two bounds in the original
\href{https://github.com/ajt60gaibb/OpenProblemsInNLA/blob/main/linear-systems-and-elimination/IE-26/README.md}{IE-26
statement}, with the exact square Fourier normalization and all its
parameter quantifiers. The conjectures and their prior analysis are due
to Austin and Trefethen, with related results in Austin's thesis and the
cited later literature. The proof is analytic; finite numerical
diagnostics are supplementary.

\textbf{Independent review.} A separate
\href{https://github.com/ajt60gaibb/OpenProblemsInNLA/blob/main/references/stepaniants-ie26-2026-09-12/verification/independent-review/IE-26-independent-review.md}{Codex-agent
full-target audit} returned PASS for both bounds, with no mathematical
correction. The
\href{https://github.com/ajt60gaibb/OpenProblemsInNLA/blob/main/references/stepaniants-ie26-2026-09-12/README.md}{submission
record} preserves the exact frozen source, independent report, and
separate coordinating check. The coordinating agent contributed before
the source was frozen and is not counted as the independent reviewer.
Substantial AI assistance is disclosed; this is informal automated
review, not external human peer review or formal verification.

\Needspace{12\baselineskip}
\subsection{1. Exact target and
conclusion}\label{exact-target-and-conclusion}

Let \(m=2N+1\), \(N\ge2\), \(h=2\pi/m\), and let \(a_j=jh\), with the
indices taken modulo \(m\). Choose arbitrary \(s_j\in[-\alpha,\alpha]\),
where \(0<\alpha<1/2\), and set \(x_j=a_j+s_jh\) modulo \(2\pi\).
Reindexing the original indices \(-N,\ldots,N\) cyclically makes no
difference. Put \(z_j=e^{ix_j}\).

Let \(\ell_j\) be the trigonometric cardinal polynomial of degree at
most \(N\) at these nodes. The Lebesgue constant is

\nopagebreak[4]
\[\Lambda_N=\sup_x\sum_j|\ell_j(x)|.\]

Let \(F=m^{-1/2}[e^{i\nu x_j}]_{j,\nu=-N}^N\) and
\(\Gamma_N=\|F^{-1}\|_2\).

\textbf{Theorem 1.} There is an absolute constant \(C\) such that

\nopagebreak[4]
\[\Lambda_N\le C\frac{N^{2\alpha}-1}{\alpha(1-2\alpha)} \qquad (N\ge2,\ 0<\alpha<1/2). \tag{1}\]

For each fixed \(1/4<\alpha<1/2\) there is \(C_\alpha\), independent of
the nodes and \(N\), such that

\nopagebreak[4]
\[\Gamma_N\le C_\alpha N^{4\alpha-1}. \tag{2}\]

Throughout the proof, \(C\) denotes an absolute positive constant,
allowed to increase from line to line. Constants marked \(C_\alpha\) may
depend on \(\alpha\). All arcs and distances are circular unless
specified otherwise. Circle integrals are with respect to ordinary
angular length; fixed factors \(2\pi\) are absorbed into absolute
constants.

\Needspace{12\baselineskip}
\newpage
\subsection{2. Interlacing and a positive-real rational
function}\label{interlacing-and-a-positive-real-rational-function}

Define

\nopagebreak[4]
\[P(z)=\prod_{j=0}^{m-1}(z-z_j),\qquad P_0(z)=z^m-1,\qquad Q(z)=z^m+1,\]

and let

\nopagebreak[4]
\[\vartheta=\frac h2\sum_j s_j,\qquad R(z)=-e^{-i\vartheta}\frac{P(z)}{Q(z)}. \tag{3}\]

The zeros \(z_j\) lie strictly between the adjacent zeros of \(Q\),
because \(|s_j|<1/2\). On the unit circle away from its poles, \(R\) is
purely imaginary. This follows directly by expressing each factor
\(e^{ix}-e^{ix_j}\) as \(2ie^{i(x+x_j)/2}\sin((x-x_j)/2)\) and using
\(Q(e^{ix})=2e^{imx/2}\cos(mx/2)\). Also

\nopagebreak[4]
\[R(0)=e^{i\vartheta},\qquad |\vartheta|\le\pi\alpha<\pi/2. \tag{4}\]

Both \(R\) and \(1/R\) have positive real part in the disk. Here is a
finite-dimensional proof, including the signs. For \(s_j=0\),
\(R=(1-z^m)/(1+z^m)\), and its partial fraction representation is a sum
of the positive-real kernels \((\eta+z)/(\eta-z)\) with coefficients
\(1/m\), at the zeros \(\eta\) of \(Q\). As the parameters change
continuously along \(\tau s_j\), \(0\le\tau\le1\), the poles stay
simple, the zeros stay strictly interlaced, and the coefficients of
these kernels remain real and nonzero, since the boundary values are
purely imaginary and no zero meets a pole. Their signs therefore remain
positive. Subtracting all pole terms leaves a rational function without
poles, hence a constant; its boundary real part is zero, so the constant
is purely imaginary. The same argument applies to \(1/R\), whose poles
are the \(z_j\). Thus, for some real \(b\),

\nopagebreak[4]
\[\frac1{R(z)}=ib+\sum_j\beta_j\frac{z_j+z}{z_j-z},\qquad \beta_j>0. \tag{5}\]

At \(z=0\), (4) gives

\nopagebreak[4]
\[\sum_j\beta_j=\cos\vartheta\le1. \tag{6}\]

Taking the magnitude of the residue at \(z_j\) gives the exact identity

\nopagebreak[4]
\[\beta_j=\frac{|Q(z_j)|}{2|P'(z_j)|}
       =\frac{\cos(\pi s_j)}{|P'(z_j)|}. \tag{7}\]

In particular,

\nopagebreak[4]
\[\frac1{|P'(z_j)|}\le\frac{\beta_j}{\cos(\pi\alpha)}. \tag{8}\]

The real part of (5) is the exact Poisson representation

\nopagebreak[4]
\[\operatorname{Re}\frac1{R(re^{ix})}
 =\sum_j\beta_j\frac{1-r^2}{|re^{ix}-z_j|^2}. \tag{9}\]

There are no zeros or poles of \(R\) in the open disk, so every
expression used here is defined.

\Needspace{12\baselineskip}
\subsection{3. A radial comparison with the exact
exponent}\label{a-radial-comparison-with-the-exact-exponent}

Put \(t_0=1/m\). For \(t\ge t_0\), write \(R_t(x)=R(e^{-t}e^{ix})\).

\textbf{Lemma 2.} Uniformly in all the parameters,

\nopagebreak[4]
\[\frac{|R_t(x)|}{|R_T(x)|}\le C(T/t)^{2\alpha}
       \qquad(t_0\le t\le T\le1), \tag{10}\]

and

\nopagebreak[4]
\[|R_t(x)|\le C t^{-2\alpha}\qquad(t_0\le t\le1). \tag{11}\]

\textbf{Proof.} Compare \(P\) with \(P_0\). The factors \(P_0/Q\) at
radii \(e^{-t}\) with \(mt\ge1\) have modulus between \(\tanh(1/2)\) and
\(\coth(1/2)\). These bounds are absolute.

For a root at angle \(\theta\),

\nopagebreak[4]
\[\frac{\partial}{\partial\theta}\log|e^{-t}e^{ix}-e^{i\theta}|
 =-K_t(x-\theta),\qquad
 K_t(u)=\frac{\sin u}{2(\cosh t-\cos u)}. \tag{12}\]

Consequently the difference between the logarithmic perturbation
products at heights \(t\) and \(T\) is bounded in absolute value by

\nopagebreak[4]
\[\alpha h\sum_j\sup_{|v-a_j|\le h/2}|K_t(x-v)-K_T(x-v)|. \tag{13}\]

For an absolutely continuous periodic function \(f\), the elementary
cell estimate

\nopagebreak[4]
\[h\sum_j\sup_{|v-a_j|\le h/2}|f(v)|
 \le\int_{-\pi}^{\pi}|f(v)|\,dv+h\int_{-\pi}^{\pi}|f'(v)|\,dv \tag{14}\]

follows by comparing a value in each cell to its average and integrating
the derivative. For \(0<t\le1\), direct differentiation of (12), or its
single positive maximum on \((0,\pi)\), gives

\nopagebreak[4]
\[\int_{-\pi}^{\pi}|K_t'(v)|\,dv\le C/t.\]

For \(0<t\le T\), \(K_t-K_T\) has the sign of \(\sin u\), and direct
integration gives

\nopagebreak[4]
\[\int_{-\pi}^{\pi}|K_t-K_T|\,du
 =2\log\frac{\coth(t/2)}{\coth(T/2)}
 \le2\log(T/t). \tag{15}\]

The last inequality uses that \(v\coth(v/2)\) is increasing. Since
\(h/t\le2\pi\), equations (13)--(15) yield \(2\alpha\log(T/t)+C\) as an
upper bound. This proves (10), after including the bounded unperturbed
factors. Letting \(T\to\infty\) in the same calculation gives
\(2\alpha\log\coth(t/2)+C\le2\alpha\log(C/t)+C\). At the origin
\(|R(0)|=1\), proving (11). The constant stays absolute because
\(2\alpha<1\). \ensuremath{\blacksquare}

\Needspace{12\baselineskip}
\subsection{4. Proof of the full uniform Lebesgue
estimate}\label{proof-of-the-full-uniform-lebesgue-estimate}

The algebraic cardinal polynomial \(P(z)/((z-z_j)P'(z_j))\) differs on
the circle from \(\ell_j\) only by a unimodular factor. For
\(|\zeta|=1\) and \(0<r<1\),

\nopagebreak[4]
\[|\zeta-z_l|^2\le r^{-1}|r\zeta-z_l|^2.\]

Applying this to the \(m-1\) factors with \(l\ne j\), with
\(r=e^{-t_0}\), gives

\nopagebreak[4]
\[|\ell_j(x)|
 \le C\frac{|P(re^{ix})|}{|re^{ix}-z_j|\,|P'(z_j)|}
 \le\frac{C}{\cos(\pi\alpha)}
       |R_{t_0}(x)|\frac{\beta_j}{|re^{ix}-z_j|}. \tag{16}\]

Here \(r^{-(m-1)/2}\le e^{1/2}\) and \(|Q(re^{ix})|\le2\). This estimate
includes evaluation exactly at a node; no division by a vanishing
boundary factor occurs.

Divide the nodes into dyadic angular-distance annuli about \(x\),
starting at size \(t_0\) and ending at size comparable to one; nodes at
distance at least one form the last set. For an annulus of size \(T\),
\(t_0\le T\le1\), its denominators in (16) are bounded below by \(cT\),
with the first annulus interpreted using the radial distance \(t_0\).
The Poisson kernel in (9), evaluated at radius \(e^{-T}\), is at least
\(c/T\) on that annulus. Thus

\nopagebreak[4]
\[\sum_{j\text{ in annulus}}\frac{\beta_j}{|re^{ix}-z_j|}
 \le C\operatorname{Re}\frac1{R_T(x)}
 \le \frac C{|R_T(x)|}. \tag{17}\]

Equations (10), (16), and (17) bound this annulus's contribution by

\nopagebreak[4]
\[\frac{C}{\cos(\pi\alpha)}(T/t_0)^{2\alpha}. \tag{18}\]

The last set contributes at most \(C|R_{t_0}(x)|/\cos(\pi\alpha)\) by
(6), and (11) gives the same bound at the largest scale. Therefore

\nopagebreak[4]
\[\sum_j|\ell_j(x)|\le
 \frac C{\cos(\pi\alpha)}\sum_{q=0}^{\lceil\log_2m\rceil}2^{2\alpha q}
 \le \frac C{\cos(\pi\alpha)}\frac{m^{2\alpha}-1}{\alpha}. \tag{19}\]

The final geometric-sum inequality is uniform for \(0<\alpha<1/2\); use
\(2^{2\alpha}-1\ge2\alpha\log2\) and absorb the bounded factor from the
last dyadic scale. Concavity of cosine on \([0,\pi/2]\) gives
\(\cos(\pi\alpha)\ge1-2\alpha\). Finally, \(m=2N+1\le(5/2)N\) and
\(N\ge2\) imply

\nopagebreak[4]
\[m^a-1\le C(N^a-1)\qquad(0<a\le1),\]

with absolute \(C\), for example by integrating \(a u^{a-1}\) and using
\(\log N\ge\log2\). Taking the supremum over \(x\) proves (1).

\Needspace{12\baselineskip}
\subsection{5. A sector function for the regularized
product}\label{a-sector-function-for-the-regularized-product}

Only the second bound remains. Its constants may depend on fixed
\(\alpha>1/4\).

Let \(s(\theta)=s_j\) on the cell \([a_j-h/2,a_j+h/2)\), periodically,
and put

\nopagebreak[4]
\[\phi(\theta)=\pi s(\theta),\qquad
 G(z)=\frac1{2\pi}\int_{-\pi}^{\pi}
       \phi(\theta)\frac{e^{i\theta}+z}{e^{i\theta}-z}\,d\theta,
 \qquad B(z)=e^{iG(z)}. \tag{20}\]

The real part of \(G\) is the Poisson integral of \(\phi\), so

\nopagebreak[4]
\[|\operatorname{Re}G(z)|\le\pi\alpha. \tag{21}\]

Thus \(B\) has an analytic logarithm with imaginary part in
\([-\pi\alpha,\pi\alpha]\), and both \(B^{1/(2\alpha)}\) and
\(B^{-1/(2\alpha)}\) have nonnegative real part.

\textbf{Lemma 3.} At \(r=e^{-1/m}\), uniformly for all angles \(x\),

\nopagebreak[4]
\[C^{-1}|B(re^{ix})|\le|R(re^{ix})|\le C|B(re^{ix})|. \tag{22}\]

Moreover, if \(|x-y|\le Ch\), with a fixed absolute multiple \(C\), then

\nopagebreak[4]
\[C_1^{-1}\le\frac{|B(re^{ix})|}{|B(re^{iy})|}\le C_1, \tag{23}\]

where \(C_1\) depends only on that fixed multiple, and is independent of
\(m\), the shifts, and \(0<\alpha<1/2\).

\textbf{Proof.} Let \(f_z(\theta)=ze^{-i\theta}/(1-ze^{-i\theta})\).
Choose all logarithms of \(1-ze^{-i\theta}\) to vanish at \(z=0\). Apart
from a unimodular constant, the logarithm of \(P/P_0\) is

\nopagebreak[4]
\[L(z)=\sum_j\{\log(1-ze^{-i(a_j+s_jh)})-\log(1-ze^{-ia_j})\}.\]

Its linear term in the node shifts is \(ih\sum_j s_j f_z(a_j)\). On the
other hand,

\nopagebreak[4]
\[G(z)-G(0)=\sum_j s_j\int_{a_j-h/2}^{a_j+h/2} f_z(\theta)\,d\theta.\]

Taylor's formula and cell quadrature therefore show

\nopagebreak[4]
\[|L(z)-i(G(z)-G(0))|
 \le C h^2\sum_j\sup_{|v-a_j|\le h/2}|f_z'(v)|\le C
       \qquad(|z|=e^{-1/m}). \tag{24}\]

For clarity, the last bound follows from

\nopagebreak[4]
\[|f_z'(v)|=\frac{|z|}{|1-ze^{-iv}|^2}
 \le \frac C{t_0^2+\operatorname{dist}(v,\arg z)^2},\]

whose cell suprema sum to at most \(C(t_0^{-2}+(ht_0)^{-1})\);
\(h/t_0=2\pi\). The factors \(\alpha\) and \(\alpha^2\) from quadrature
and Taylor are bounded absolutely. Since \(G(0)\) is real and \(P_0/Q\)
is bounded above and below at this radius, (24) proves (22).

Differentiating the kernel in (20) gives

\nopagebreak[4]
\[|\partial_xG(re^{ix})|\le C\alpha/t_0.\]

Integration over a distance \(O(h)\) proves (23). The same derivative
estimate applies if \(\phi\) is restricted to any measurable arc. \ensuremath{\blacksquare}

\Needspace{12\baselineskip}
\subsection{6. Local second moments of a sector
function}\label{local-second-moments-of-a-sector-function}

Set

\nopagebreak[4]
\[w_j=|B(re^{ia_j})|,\qquad r=e^{-1/m}.\]

\textbf{Lemma 4.} Fix \(1/4<\alpha<1/2\). For every circular interval
\(I\) of \(q\) consecutive grid indices, \(1\le q\le m\), there is a
positive number \(b_I\) such that

\nopagebreak[4]
\[\sum_{j\in I}w_j^2\le C_\alpha b_I^2q^{4\alpha},\qquad
 \sum_{j\in I}w_j^{-2}\le C_\alpha b_I^{-2}q^{4\alpha}. \tag{25}\]

\textbf{Proof.} We first specify the one standard harmonic-analysis
estimate used here. If \(f\) is analytic in a neighborhood of the closed
disk and \(\operatorname{Re}f\ge0\) in the disk, then

\nopagebreak[4]
\[|\{\theta:|f(e^{i\theta})|>u\}|
       \le C|f(0)|/u. \tag{26}\]

Indeed, its real part has integral \(2\pi\operatorname{Re}f(0)\), while
the imaginary part minus \(\operatorname{Im}f(0)\) is the periodic
Hilbert transform of its real part. Markov's inequality, the weak
\((1,1)\) inequality for that Hilbert transform, and a separate bound
for the constant imaginary part give (26). The precise periodic weak
theorem is Theorem 12.1 of Laugesen {[}4{]}, applicable here to smooth,
hence \(L^2\), boundary data. Composition with a disk automorphism gives
the corresponding harmonic-measure estimate with any interior point in
place of zero.

Suppose first that \(\ell=qh\le1/100\). Choose an arc \(J\), centered at
the middle of \(I\), of length \(8\ell\), and split
\(G=G_{\rm loc}+G_{\rm far}\) using \(\phi\mathbf1_J\) and its
complement in (20). On the arc containing \(I\) and an extra cell on
either side, at height \(t_0\), the variation of \(G_{\rm far}\) is at
most \(C\alpha\): the derivative of its kernel is bounded by
\(C/\operatorname{dist}^2\), whose integral outside \(J\) is at most
\(C/\ell\), and the relevant angular variation is \(O(\ell)\). Set

\nopagebreak[4]
\[b_I=\exp(-\operatorname{Im}G_{\rm far}(re^{ix_I})),\qquad
 B_{\rm loc}=e^{iG_{\rm loc}},\]

where \(x_I\) is the center angle. Thus \(w_j/b_I\) and
\(|B_{\rm loc}(re^{ia_j})|\) are comparable by absolute factors, and the
same holds for their reciprocals.

Put \(p=1/(2\alpha)\in(1,2)\). Both functions

\nopagebreak[4]
\[f_\pm(z)=\exp(\pm iG_{\rm loc}(z)/(2\alpha))\]

have nonnegative real part. Let \(z_I=e^{-\ell}e^{ix_I}\). Since
\(\phi\mathbf1_J\) is supported on an arc of length \(8\ell\) and the
Herglotz kernel at \(z_I\) has magnitude at most \(C/\ell\) there,

\nopagebreak[4]
\[|G_{\rm loc}(z_I)|\le C\alpha,\qquad |f_\pm(z_I)|\le C. \tag{27}\]

Apply (26) to \(f_\pm(rz)\) after a disk automorphism sending zero to
\(e^{-(\ell-t_0)}e^{ix_I}\). This is legitimate because
\(\ell=qh\ge2\pi t_0>2t_0\). The resulting harmonic-measure density on
the angular arc containing \(I\) and its neighboring cells is bounded
below by \(c/\ell\). Consequently,

\nopagebreak[4]
\[|\{x\text{ in that arc}:|B_{\rm loc}(re^{ix})|^{\pm1}>u\}|
       \le C\ell u^{-p}. \tag{28}\]

The derivative estimate in Lemma 3 shows that \(|B_{\rm loc}|\) and its
inverse vary by at most an absolute factor on each grid cell. Applying
(28) to the disjoint cells around the selected grid points gives

\nopagebreak[4]
\[\#\{j\in I:|B_{\rm loc}(re^{ia_j})|^{\pm1}>u\}
       \le Cq u^{-p}. \tag{29}\]

There is also the sharp pointwise cutoff

\nopagebreak[4]
\[|B_{\rm loc}(re^{ix})|^{\pm1}\le Cq^{2\alpha} \tag{30}\]

on these cells. To verify its exponent, the imaginary part of the kernel
in (20) is

\nopagebreak[4]
\[\frac{\sin(x-\theta)}{\cosh t_0-\cos(x-\theta)}.\]

Since \(|\phi|\le\pi\alpha\), integration of its absolute value over the
support \(J\), which lies within distance \(C\ell\) of \(x\), yields

\nopagebreak[4]
\[|\operatorname{Im}G_{\rm loc}(re^{ix})|
 \le 2\alpha\log(\ell/t_0)+C\alpha.
\]

This follows by the elementary antiderivative \(\log(\cosh t_0-\cos u)\)
on each half-interval; \(\ell/t_0=2\pi q\). It proves (30).

For either sign, integrate the counting function in (29), using (30) and
the trivial count \(q\) below level one:

\nopagebreak[4]
\[\sum_{j\in I}|B_{\rm loc}(re^{ia_j})|^{\pm2}
 \le Cq+Cq\int_1^{Cq^{2\alpha}}u^{1-p}\,du
 \le C_\alpha q^{4\alpha}. \tag{31}\]

Here \(2-p>0\) is precisely the strict assumption \(\alpha>1/4\).
Restoring \(b_I\) proves (25) for small arcs.

If \(qh>1/100\), use the whole circle as the local region and take
\(b_I=1\). Now \(G(0)\) is real, so \(|\exp(\pm iG(0)/(2\alpha))|=1\).
Apply (26) directly to \(\exp(\pm iG(rz)/(2\alpha))\), use cell
comparability, and integrate the full-circle kernel to get the cutoff
\(Cm^{2\alpha}\). The same calculation proves the bounds with
\(m^{4\alpha}\). Since \(m\le Cq\) in this case and \(4\alpha<2\), these
imply (25), with absolute comparison factors. \ensuremath{\blacksquare}

\Needspace{12\baselineskip}
\subsection{7. Cardinal matrix and the no-logarithm spectral
estimate}\label{cardinal-matrix-and-the-no-logarithm-spectral-estimate}

Let \(E\) be the interpolation matrix from the perturbed nodes to the
equispaced nodes:

\nopagebreak[4]
\[E_{kj}=\ell_j(a_k).\]

If \(F_0=m^{-1/2}[e^{i\nu a_k}]\) is the unitary Fourier matrix, then

\nopagebreak[4]
\[E=F_0F^{-1},\qquad \|E\|_2=\Gamma_N. \tag{32}\]

We claim the pointwise bound

\nopagebreak[4]
\[|E_{kj}|\le C_\alpha\frac{w_k/w_j}{1+d(k,j)},\qquad
 d(k,j)=\min_{v\in\mathbb Z}|k-j+vm|. \tag{33}\]

To prove it, (9) at \(rz_j\) gives

\nopagebreak[4]
\[\beta_j\frac{1+r}{1-r}\le\operatorname{Re}\frac1{R(rz_j)}
 \le\frac1{|R(rz_j)|},\]

so \(\beta_j\le C m^{-1}|R(rz_j)|^{-1}\). By (7), (22), and (23),

\nopagebreak[4]
\[\frac1{|P'(z_j)|}\le C_\alpha m^{-1}w_j^{-1}.\]

The product inequality preceding (16) gives

\nopagebreak[4]
\[|E_{kj}|\le C\frac{|P(re^{ia_k})|}{|re^{ia_k}-z_j|\,|P'(z_j)|}.\]

Lemma 3 and the bounded \(Q\) factor bound the numerator by \(Cw_k\).
Finally,

\nopagebreak[4]
\[|re^{ia_k}-z_j|\ge c\frac{1+d(k,j)}m,\]

because \(t_0=1/m\) handles \(d=0\), and for \(d\ge1\) the angular
distance is at least \((d-\alpha)h\ge dh/2\). This proves (33).

Decompose \(E\) according to the dyadic distance ranges \(d<2\), and
\(R\le d<2R\) with \(R=2,4,8,\ldots\), up to the circle's diameter. At a
fixed scale \(R\), partition the cyclic indices into consecutive blocks
of length at most \(R\). Each block interacts in this distance range
with only an absolute number of other blocks. The union of two
interacting blocks lies in a circular interval containing at most \(CR\)
indices, or is covered by the whole circle when \(R\) is comparable to
\(m\).

For an interacting row block \(I_1\) and column block \(I_2\), (33) and
Cauchy--Schwarz give the following norm bound for the block, including
any mask selecting the distance range:

\nopagebreak[4]
\[\|E^{(R)}_{I_1,I_2}\|_2
 \le\frac{C_\alpha}{R}
       \Big(\sum_{k\in I_1}w_k^2\Big)^{1/2}
       \Big(\sum_{j\in I_2}w_j^{-2}\Big)^{1/2}
 \le C_\alpha R^{4\alpha-1}. \tag{34}\]

The last step applies Lemma 4 to one interval containing both blocks, so
its two normalization factors cancel. The same estimate, with a
constant, handles \(d<2\). A block matrix with uniformly bounded numbers
of nonzero blocks in every block row and column has norm at most an
absolute multiple of its largest block norm: apply the scalar row/column
bound to the matrix of block norms. Hence (34) gives

\nopagebreak[4]
\[\|E^{(R)}\|_2\le C_\alpha R^{4\alpha-1}.\]

Since \(4\alpha-1>0\), summing the dyadic scales is a geometric sum:

\nopagebreak[4]
\[\Gamma_N=\|E\|_2
 \le C_\alpha\sum_{R\text{ dyadic},\ R\le m}R^{4\alpha-1}
 \le C_\alpha m^{4\alpha-1}
 \le C_\alpha N^{4\alpha-1}. \tag{35}\]

This proves (2). Together with Section 4 it proves Theorem 1 for the
full original bundled target. No assertion at \(\alpha=1/4\) is made; at
that endpoint the second-moment integral in (31) and the scale summation
require different estimates.

\Needspace{12\baselineskip}
\subsection{8. Attribution, scope, and
verification}\label{attribution-scope-and-verification}

The two exact targets originate in Austin and Trefethen {[}1{]} and
Austin's thesis {[}2{]}. Chen, Lin, and Zhang {[}3{]} prove the second
bound with an additional logarithm; the present argument removes it. The
only imported analytic inequality beyond elementary complex analysis is
the periodic Hilbert transform's weak \((1,1)\) estimate, in Laugesen
{[}4{]}, Theorem 12.1, printed p.67. Its applicability to the smooth
boundary data used in Lemma 4 is explained there. All grid-product,
residue, localization, and matrix estimates needed for this conclusion
appear in Sections 1--7.

Nodes index the rows of the normalized Fourier matrix in this proof,
exactly as in the canonical statement. A convention with frequency
indices as rows transposes this matrix and preserves all singular
values, including the inverse norm. The identity \(E=F_0F^{-1}\) in (32)
uses the displayed row convention.

The first estimate has one absolute constant simultaneously for all
\(0<\alpha<1/2\). The second constant may depend on a fixed
\(1/4<\alpha<1/2\). No estimate at \(\alpha=1/4\), no assertion about
optimal constants, and no rectangular oversampling or
restricted-isometry statement is made.

The original frozen source, dated public eligibility check, numerical
identity checks, and review reports are retained in the
\href{https://github.com/ajt60gaibb/OpenProblemsInNLA/blob/main/references/stepaniants-ie26-2026-09-12/README.md}{submission
record}. The diagnostics test three deterministic instances and do not
establish any quantified theorem. The mathematical proof is independent
of those tests. Source and eligibility checks are bounded; no priority
claim is made.

\Needspace{12\baselineskip}
\subsection{References}\label{references}

\begin{enumerate}
\def\labelenumi{\arabic{enumi}.}
\tightlist
\item
  A. P. Austin and L. N. Trefethen, \emph{Trigonometric interpolation
  and quadrature in perturbed points}, SIAM J. Numer. Anal. 55 (2017),
  2113--2122. The conjectured replacement in equation (12),
  pp.2115--2116, and the second-norm conjecture, p.2119.
  \href{https://doi.org/10.1137/16M1107760}{DOI};
  \href{https://people.maths.ox.ac.uk/trefethen/perturbed.pdf}{author
  PDF}.
\item
  A. P. Austin, \emph{Some New Results on and Applications of
  Interpolation in Numerical Computation}, D.Phil. thesis, University of
  Oxford, 2016. Conjectures 3.5 and 3.10, printed pp.75 and 82.
  \href{https://personal.math.vt.edu/apaustin/pubs/DPhilThesis.pdf}{Author
  PDF}. The endpoint assertion of the thesis is not imported into (2).
\item
  L. Chen, R. Lin, and H. Zhang, \emph{The smallest singular value of
  nonuniform Fourier matrices}, arXiv:2608.21960v1, 22 August 2026.
  Theorems 1.4--1.5 and Section 3.
  \href{https://arxiv.org/abs/2608.21960v1}{Versioned record};
  \href{https://arxiv.org/html/2608.21960v1}{full text}.
\item
  R. S. Laugesen, \emph{Harmonic Analysis Lecture Notes},
  arXiv:0903.3845v2, 17 May 2017, Theorem 12.1, printed p.67.
  \href{https://arxiv.org/abs/0903.3845v2}{Versioned record};
  \href{https://arxiv.org/pdf/0903.3845v2}{PDF}.
\end{enumerate}
\end{document}
